\documentclass[11pt]{article}
\usepackage{amsmath,amssymb,amsfonts}
\usepackage[left=2.5cm,right=2.5cm,top=2cm,bottom=2cm]{geometry}
\usepackage{booktabs}
\usepackage{hyperref}

\DeclareMathOperator{\Var}{Var}
\DeclareMathOperator{\Cov}{Cov}
\DeclareMathOperator{\Corr}{Corr}
\DeclareMathOperator{\tr}{tr}
\newcommand{\bI}{\mathbf{I}}
\newcommand{\bone}{\mathbf{1}}
\newcommand{\bw}{\mathbf{w}}

\title{Notebook Extensions:\\
  General Time Points, AR(1) Cross-Covariance,\\
  and the $K$-Visit Conventional Design}
\author{R.G.\ Thomas}
\date{Transcribed from handwritten notes, Jul 2019 -- Nov 2021}

\begin{document}
\maketitle
\tableofcontents


\section{Frost 3.1.1 with general time points}
\label{sec:general_time}

\noindent
\textit{(From notes dated 11/16, pages 1--6.)}

\medskip
\noindent
Rederive the Frost conventional-design variance keeping
general time points $t_1, t_2$ (not equally spaced), subject
to $\sum t_i = 0$ and $\sum t_i^2 = 1$.

\subsection{Setup}

Change scores $C_1 = Y_1 - Y_0$, $C_2 = Y_2 - Y_0$.
Design:
\begin{equation}
  X\beta + Zb
  = \begin{pmatrix} t_1 & t_1 q \\ t_2 & t_2 q \end{pmatrix}
    \begin{pmatrix} \beta \\ \gamma \end{pmatrix}
  + \begin{pmatrix} t_1 \\ t_2 \end{pmatrix} b_i
\end{equation}

Residual covariance and inverse:
\begin{equation}
  R = a\begin{pmatrix} 2 & 1 \\ 1 & 2 \end{pmatrix},
  \qquad
  R^{-1} = \frac{1}{3a}\begin{pmatrix}
    2 & -1 \\ -1 & 2
  \end{pmatrix}
\end{equation}

\subsection{$X'R^{-1}X$ with general $t_1, t_2$}

For treatment ($q = g_i$):
\begin{align}
  X'R^{-1}X
  &= \frac{1}{3a}
    \begin{pmatrix} t_1 & t_2 \\ t_1 q & t_2 q \end{pmatrix}
    \begin{pmatrix} 2&-1\\-1&2 \end{pmatrix}
    \begin{pmatrix} t_1 & t_1 q \\ t_2 & t_2 q \end{pmatrix}
    \notag \\[4pt]
  &= \frac{1}{3a}\begin{pmatrix}
    2(t_1^2 + t_2^2) - 2t_1 t_2 &
    2q(t_1^2 + t_2^2) - 2q t_1 t_2 \\[2pt]
    2q(t_1^2 + t_2^2) - 2q t_1 t_2 &
    2q^2(t_1^2 + t_2^2) - 2q^2 t_1 t_2
  \end{pmatrix}
\end{align}

Using $t_1^2 + t_2^2 = 1$:
\begin{equation}
  X'R^{-1}X = \frac{2(1 - t_1 t_2)}{3a}
  \begin{pmatrix} 1 & q \\ q & q^2 \end{pmatrix}
\end{equation}


\subsection{$X'R^{-1}Z$ and the Woodbury scalar $H$}

\begin{equation}
  X'R^{-1}Z = \frac{1}{3a}
    \begin{pmatrix} t_1 & t_2 \\ t_1 q & t_2 q \end{pmatrix}
    \begin{pmatrix} 2&-1\\-1&2 \end{pmatrix}
    \begin{pmatrix} t_1 \\ t_2 \end{pmatrix}
  = \frac{2(1 - t_1 t_2)}{3a}
    \begin{pmatrix} 1 \\ q \end{pmatrix}
\end{equation}

\begin{equation}
  Z'R^{-1}Z = \frac{2(1 - t_1 t_2)}{3a}
\end{equation}

\begin{equation}
  H = (G^{-1} + Z'R^{-1}Z)^{-1}
  = \left(\frac{1}{b}
    + \frac{2(1 - t_1 t_2)}{3a}\right)^{-1}
  = \frac{3ab}{3a + 2b(1 - t_1 t_2)}
  \label{eq:H_general}
\end{equation}


\subsection{Information matrix and variance}

After the Woodbury correction (summing treatment and
placebo):
\begin{align}
  X'\Sigma^{-1}X
  &= \frac{2(1 - t_1 t_2)}{3a}
    \begin{pmatrix} 2 & 1 \\ 1 & 1 \end{pmatrix}
    - \frac{2b(1-t_1 t_2)}{3a(3a + 2b(1-t_1 t_2))}
    \begin{pmatrix} 2 & 1 \\ 1 & 1 \end{pmatrix}
    \notag \\[4pt]
  &= \frac{2(1-t_1 t_2)}{3a + 2b(1-t_1 t_2)}
    \begin{pmatrix} 2 & 1 \\ 1 & 1 \end{pmatrix}
    \cdot \frac{1}{1}
\end{align}

Inverting:
\begin{equation}
  (X'\Sigma^{-1}X)^{-1}
  = \frac{3a + 2b(1 - t_1 t_2)}{2(1 - t_1 t_2)}
    \begin{pmatrix} 1 & -1 \\ -1 & 2 \end{pmatrix}
\end{equation}

The $(2,2)$ element gives:
\begin{equation}
  \boxed{
    \Var(\hat\gamma)
    = \frac{3a + 2b(1 - t_1 t_2)}{1 - t_1 t_2}
    = \frac{3a}{1 - t_1 t_2} + 2b
  }
  \label{eq:var_general}
\end{equation}


\subsection{Verification with equally spaced times}

\textbf{Case $t_1 = t$, $t_2 = 2t$:}
$1 - t_1 t_2 = 1 - 2t^2$. However, the normalization
$t_1^2 + t_2^2 = 1$ gives $5t^2 = 1$, so
$1 - 2t^2 = 1 - 2/5 = 3/5$. Then
$\Var(\hat\gamma) = 5a + 2b$.

\textbf{Case $t_1 + t_2 = 0$} (centered):
$1 - t_1 t_2 = 1 - t_1(-t_1) = 1 + t_1^2$.

\textbf{Case $t_1 = 1$, $t_2 = 2$} (unnormalized,
matching Frost):
$1 - t_1 t_2 = -1$. Then $3a/(-1) + 2b = -3a + 2b$
is negative, indicating the normalization $t_1^2+t_2^2=1$
is essential. With Frost's parameterization (no normalization,
$R = a\begin{pmatrix}2&-1\\-1&2\end{pmatrix}$), the
circled result on page 3 is
$\frac{4b}{a(a+2b)}$, which is the per-pair variance
in Frost's original units.


\subsection{Correlation form}

With $c = a + 2b$ and $r = (b - a)/(b + 2a)$:
\begin{equation}
  1 - r^2 = \frac{(b+2a)^2 - (b-a)^2}{(b+2a)^2}
  = \frac{3a^2 + 6ab}{(b+2a)^2}
  = \frac{3a(a + 2b)}{(b+2a)^2}
  = \frac{3a \cdot c}{(b+2a)^2}
\end{equation}

Alternative parameterization explored (page 15):
$c = 2a + b$, $d = a + 2b$. Then
$a = (c - b)/2$, $b = (2d - c)/3$.
Circled: $\frac{2d - c}{d}$.


\section{AR(1) cross-covariance between run-in and treatment}
\label{sec:ar1_cross}

\noindent
\textit{(From notes dated 10/13, pages 7--12.)}

\medskip
\noindent
Under an AR(1) correlation structure with parameter
$c = \rho$, the cross-covariance between the run-in mean
(over $r$ observations) and the treatment-period mean
(over $q$ observations) separated by a gap of $K$ time
points is:
\begin{equation}
  \Cov(\bar{X}_0, \bar{X}_K)
  = \frac{\sigma^2}{rc}\,\rho_1\, c^{K-1}\,
    \frac{(1 - c^r)(1 - c^q)}{(1 - c)^2}
  \label{eq:cov_ar1}
\end{equation}

\subsection{Derivation}

The covariance of the means involves a double sum:
\begin{equation}
  \Cov(\bar{X}_0, \bar{X}_K)
  = \frac{1}{rc}\sum_{i=1}^{r}\sum_{j=1}^{c}
    \Cov(R_i, C_j)
\end{equation}

Under AR(1), $\Cov(R_i, C_j) = \sigma^2 \rho^{|k_j - k_i|}$
where $k_j - k_i$ is the lag. With $K$ gap between the
last run-in and first treatment visit:
\begin{equation}
  \Cov(R_i, C_j) = \sigma^2 \rho^{K + (r - i) + (j - 1)}
  = \sigma^2 \rho^{K-1} \cdot \rho^{r-i} \cdot \rho^{j}
\end{equation}

The double sum factors:
\begin{equation}
  \frac{\sigma^2}{rc}\, \rho^{K-1}
    \sum_{i=1}^{r} \rho^{r-i}
    \sum_{j=1}^{q} \rho^{j}
  = \frac{\sigma^2}{rc}\, \rho^{K-1}\,
    \frac{1 - \rho^r}{1 - \rho}\,
    \frac{\rho(1 - \rho^q)}{1 - \rho}
\end{equation}

which gives \eqref{eq:cov_ar1} with $\rho_1 = \rho$.

\subsection{Closed form for the double sum}

The key identity (referencing Knuth, p.\ 40) is the
evaluation of
\begin{equation}
  \sum_{j=0}^{q}\;\sum_{i=j}^{r-1+j} c^i
  = \sum_{j=0}^{q}\; c^j \sum_{i=0}^{r-1} c^i
  = \frac{(1 - c^{r+1})(1 - c^{q+1})}{c(1-c)^2}
  \label{eq:double_sum}
\end{equation}

Verification for $r = 3$, $q = 1$:
\begin{equation}
  \frac{(1 - c^4)(1 - c^2)}{c(1-c)^2}
  = \frac{(1-c^2)(1+c^2)(1-c)(1+c)}{c(1-c)^2}
  = \frac{(1+c)(1+c^2)}{c} \cdot \frac{1-c^2}{1-c}
\end{equation}

Cross-checked by direct expansion:
$1 + c + c^2 + c^3 + c^4$ summed appropriately.


\section{Stratified variance with unequal visit counts}
\label{sec:stratified}

\noindent
\textit{(From notes dated 10/13, pages 13--14.)}

\medskip
\noindent
When participants have varying numbers of observations,
the variance of the difference in group means is:
\begin{equation}
  \Var(\bar{X}_K - X_0)
  = V(X_K) + V(X_0) - 2\Cov(X_K, X_0)
\end{equation}

With change scores $z_i = X_{Ki} - X_{0i}$ and
$w_i = Y_{Ki} - Y_{0i}$:
\begin{equation}
  \Var(\bar{z} - \bar{w})
  = \frac{1}{n^2}\Var\!\left(\sum z_i\right)
    + \frac{1}{m^2}\Var\!\left(\sum w_i\right)
\end{equation}

Expanding within strata indexed by number of visits
$n_s$ (with $s$ strata, example weights 40\%, 30\%,
20\%, 10\%):
\begin{equation}
  \frac{1}{n^2}\left[
    m_1\,\Var(X_K - X_0)\big|_{n=5}
    + 2m_2\,\Var(X_K - X_0)\big|_{n=4}
    + 3m_3\,\Var(X_K - X_0)\big|_{n=3}
    + 4m_4\,\Var(X_K - X_0)\big|_{n=2}
  \right]
\end{equation}
where $m_1 + m_2 + \ldots = n$ and the per-stratum
variance uses the AR(1) formula from
Section~\ref{sec:ar1_cross}.


\section{Conventional design with $K$ visits}
\label{sec:K_visits}

\noindent
\textit{(From unnumbered pages, TRT/PLA derivation.)}

\medskip
\noindent
For the conventional design (no run-in) with $K$
post-baseline visits, the raw-outcome Woodbury identity
gives a clean closed-form result.

\subsection{Setup}

Let $\bw = (t_1, \ldots, t_K)' = d(1, 2, \ldots, K)'$
be the time vector. The design matrices are:
\begin{align}
  X^{(\text{trt})} &= (\bw,\; \bw)
    & &\text{(columns for $\beta$ and $\gamma$)} \\
  X^{(\text{plc})} &= (\bw,\; \mathbf{0})
\end{align}

Random effects design: $U = \bw$, $G = \sigma_1^2$,
$R = \sigma^2 \bI_K$.

The Woodbury inverse is
\begin{equation}
  V^{-1} = \frac{1}{\sigma^2}\bI
    - \frac{1}{\sigma^2}\cdot
    \frac{\sigma_1^2}{\sigma^2 + (K-1)\sigma_1^2}\,
    \bw\bw'
\end{equation}

Define $\rho = \sigma_1^2 / (\sigma^2 + \sigma_1^2)$
(intraclass correlation for adjacent visits).


\subsection{Treatment group information}

\begin{equation}
  \frac{1}{\sigma^2} X'X
  = \frac{1}{\sigma^2}\begin{pmatrix}\bw\\\bw\end{pmatrix}'
    \begin{pmatrix}\bw & \bw\end{pmatrix}
  = \frac{K-1}{\sigma^2}
    \begin{pmatrix} 1 & 1 \\ 1 & 1 \end{pmatrix}
\end{equation}
(using $\bw'\bw = \sum t_i^2 = K - 1$ under
normalization).

The Woodbury correction:
\begin{equation}
  X'\bw = \begin{pmatrix} K-1 \\ K-1 \end{pmatrix},
  \qquad
  \bw'\bw \cdot X = (K-1)^2
    \begin{pmatrix} 1 & 1 \\ 1 & 1 \end{pmatrix}
\end{equation}

\begin{align}
  X_i^{(\text{trt})\prime} V^{-1} X_i^{(\text{trt})}
  &= \frac{K-1}{\sigma^2}
    \begin{pmatrix}1&1\\1&1\end{pmatrix}
    - \frac{1}{\sigma^2}\cdot
    \frac{\sigma_1^2}{\sigma^2 + (K-1)\sigma_1^2}
    (K-1)^2
    \begin{pmatrix}1&1\\1&1\end{pmatrix}
    \notag \\[4pt]
  &= \frac{K-1}{\sigma^2}\left[
    1 - \frac{(K-1)\rho}{1 + (K-1)\rho}
  \right]
  \begin{pmatrix}1&1\\1&1\end{pmatrix}
    \notag \\[4pt]
  &= \frac{K-1}{\sigma^2}\cdot
    \frac{1}{1 + (K-1)\rho}
    \begin{pmatrix}1&1\\1&1\end{pmatrix}
\end{align}


\subsection{Placebo group information}

\begin{align}
  X_i^{(\text{plc})\prime} V^{-1} X_i^{(\text{plc})}
  &= \frac{K-1}{\sigma^2}\cdot
    \frac{1}{1 + (K-1)\rho}
    \begin{pmatrix}1&0\\0&0\end{pmatrix}
\end{align}


\subsection{Total information and variance}

Summing over $m$ per group:
\begin{equation}
  X'\Sigma^{-1}X
  = \frac{m(K-1)}{\sigma^2}\cdot
    \frac{1}{1 + (K-1)\rho}
    \begin{pmatrix}2&1\\1&1\end{pmatrix}
\end{equation}

Inverting:
\begin{equation}
  (X'\Sigma^{-1}X)^{-1}
  = \frac{\sigma^2}{m(K-1)}\bigl(1 + (K-1)\rho\bigr)
    \begin{pmatrix}1&-1\\-1&2\end{pmatrix}
\end{equation}

The $(2,2)$ element:
\begin{equation}
  \boxed{
    \Var(\hat\gamma)
    = \frac{2\sigma^2}{m(K-1)}
      \bigl(1 + (K-1)\rho\bigr)
    = \frac{2}{m(K-1)}
      \bigl(\sigma^2 + (K-1)\sigma_1^2\bigr)
  }
  \label{eq:var_K}
\end{equation}

For $K = 2$ (Frost case): $\Var(\hat\gamma) =
\frac{2}{m}(\sigma^2 + \sigma_1^2)$. $\checkmark$


\section{Two run-in, one follow-up: covariance structure}
\label{sec:r2k1_R}

\noindent
\textit{(From pages 16--18.)}

\medskip
\noindent
For $r = 2$ run-in and $k = 1$ follow-up, the change
scores relative to baseline are:
\begin{equation}
  C_1 = Y_0 - Y_{-2}, \quad
  C_2 = Y_0 - Y_{-1}, \quad
  C_3 = Y_1 - Y_0
\end{equation}

Note the sign convention: $C_1$ and $C_2$ are
\emph{forward} changes (baseline minus earlier), while
$C_3$ is the usual post-baseline change. Under the
model with residual variance $\sigma^2$:
\begin{equation}
  R = a\begin{pmatrix}
    2 & 1 & -1 \\ 1 & 2 & -1 \\ -1 & -1 & 2
  \end{pmatrix}
\end{equation}

However, this is crossed out in the notes and replaced
with the standard change-from-baseline convention
($C_1 = Y_{-2} - Y_0$, $C_2 = Y_{-1} - Y_0$,
$C_3 = Y_1 - Y_0$), which gives
\begin{equation}
  R = a(\bI_3 + \bone_3 \bone_3')
  = a\begin{pmatrix}
    2 & 1 & 1 \\ 1 & 2 & 1 \\ 1 & 1 & 2
  \end{pmatrix}
\end{equation}

The inverse uses the Sherman--Morrison formula. With
$B = \bone\bone'$ ($3 \times 3$, all ones) and
$\tr(B) = 3$:
\begin{equation}
  R^{-1} = \frac{1}{a}\left(
    \bI - \frac{1}{4}\bone\bone'
  \right)
  = \frac{1}{4a}\begin{pmatrix}
    3 & -1 & 1 \\ -1 & 3 & 1 \\ 1 & 1 & 3
  \end{pmatrix}
\end{equation}

Wait --- the $(1,3)$ and $(2,3)$ entries should be
$-1$, not $+1$. The correct inverse (matching the
report) is
\begin{equation}
  R^{-1} = \frac{1}{4a}\begin{pmatrix}
    3 & -1 & -1 \\ -1 & 3 & -1 \\ -1 & -1 & 3
  \end{pmatrix}
\end{equation}
since $R = a(\bI + \bone\bone')$ implies
$R^{-1} = \frac{1}{a}(\bI - \frac{1}{p+1}\bone\bone')
= \frac{1}{a}(\bI - \frac{1}{4}\bone\bone')$.


\section{Conceptual foundations}
\label{sec:concepts}

\noindent
\textit{(From earliest notes dated 7/27/19.)}

\medskip
\noindent
``Look at Chan simulation 2018 (Lilect). Understand
variance components!''

The conceptual picture: individual trajectories diverge
from a common baseline, with between-subject slope
variability $\sigma_b^2$ and within-subject measurement
error $\sigma^2$. The key insight is that when
$\sigma_b / \sigma$ is large, pre-randomization
observations are highly informative about each
participant's underlying trajectory, and incorporating
this information reduces the variance of the treatment
effect estimator.

The full model with group-specific effects:
\begin{equation}
  Y_{ij} = \alpha + (\alpha_j + \alpha_{ij})
    + (\beta + \beta_j + b_{ij} + \gamma g_i)\,t
    + e_i
\end{equation}
where $\alpha_j$ and $\beta_j$ are group-level intercept
and slope perturbations, $\alpha_{ij}$ and $b_{ij}$ are
individual-level random effects, and $\gamma$ is the
treatment effect.


\section{Notes consolidation plan}
\label{sec:consolidation}

\noindent
\textit{(From notes dated 7/8/21.)}

\medskip
\noindent
``Wu paper. Run-in notes. Plan: notes consolidation.''

Defines the general change score incorporating run-in
and common close:
\begin{equation}
  d_{ik} = \frac{1}{c_i + 1}
    \sum_{\ell=0}^{c_i} Y_{i,q+\ell,k}
  - \frac{1}{r_i + 1}
    \sum_{\ell=0}^{r_i} Y_{i,\ell,k}
\end{equation}

Special cases:
\begin{itemize}
  \item $r_i = c_i = 0$: conventional
    $d_{ik} = Y_{iqk} - Y_{i0k}$

  \item $r_i = c_i = 1$:
    $d_{ik} = \frac{Y_{iq} + Y_{iq+1}}{2}
      - \frac{Y_{i,-1} + Y_{i,0}}{2}$
\end{itemize}

This averaging formulation is an alternative to the
GLS approach developed in the report, and corresponds
to a simple $t$-test on averaged change scores.

\end{document}
